% pracovni_list.tex -- Lekce 04: Cykly a animace
\newcommand{\lekce}{Lekce 04 -- Pracovní list}
% preamble-pl.tex -- sdílená preamble pro pracovní listy
% Includuje se z ucitel/lekceXX/pracovni_list.tex jako:
%   % preamble-pl.tex -- sdílená preamble pro pracovní listy
% Includuje se z ucitel/lekceXX/pracovni_list.tex jako:
%   % preamble-pl.tex -- sdílená preamble pro pracovní listy
% Includuje se z ucitel/lekceXX/pracovni_list.tex jako:
%   \input{../spolecne/preamble-pl.tex}

\documentclass[a4paper,12pt]{article}

% --- Jazyk a font ---
\usepackage{polyglossia}
\setmainlanguage{czech}
\usepackage{fontspec}

% --- Okraje ---
\usepackage[top=20mm, bottom=20mm, left=22mm, right=22mm]{geometry}

% --- Česky korektní uvozovky ---
\usepackage{csquotes}
\newcommand{\uv}[1]{\enquote{#1}}

% --- Typografie ---
\usepackage{microtype}
\usepackage{parskip}          % odstavce odděleny mezerou, ne odsazením

% --- Barvy a grafika ---
\usepackage[table]{xcolor}
\definecolor{microbit-blue}{HTML}{1E4D8C}
\definecolor{code-bg}{HTML}{F5F5F5}
\definecolor{task-bg}{HTML}{EEF4FF}

% --- Tabulky ---
\usepackage{booktabs}
\usepackage{multirow}

% --- Zdrojový kód (Python) ---
\usepackage{listings}
\lstset{
  language=Python,
  basicstyle=\ttfamily\small,
  backgroundcolor=\color{code-bg},
  frame=single,
  framerule=0pt,
  rulecolor=\color{gray!30},
  xleftmargin=6pt,
  xrightmargin=6pt,
  breaklines=true,
  showstringspaces=false,
  keywordstyle=\color{microbit-blue}\bfseries,
  stringstyle=\color{teal},
  commentstyle=\color{gray}\itshape,
  literate=
    {á}{{\'a}}1 {é}{{\'e}}1 {í}{{\'i}}1 {ó}{{\'o}}1 {ú}{{\'u}}1
    {ů}{{\r{u}}}1 {č}{{\v{c}}}1 {š}{{\v{s}}}1 {ž}{{\v{z}}}1
    {Á}{{\'A}}1 {É}{{\'E}}1 {Í}{{\'I}}1 {Ó}{{\'O}}1 {Ú}{{\'U}}1
    {Č}{{\v{C}}}1 {Š}{{\v{S}}}1 {Ž}{{\v{Z}}}1 {ň}{{\v{n}}}1 {ř}{{\v{r}}}1,
}

% --- Zadání úkolů ---
\usepackage[most]{tcolorbox}
\tcbuselibrary{breakable}

\newcounter{ukol}[section]
\newenvironment{ukol}[1][]{%
  \stepcounter{ukol}%
  \begin{tcolorbox}[
    colback=task-bg,
    colframe=microbit-blue,
    fonttitle=\bfseries,
    title={Úkol~\theukol\ifx\relax#1\relax\else{: #1}\fi},
    breakable,
    left=6pt, right=6pt, top=4pt, bottom=4pt
  ]%
}{%
  \end{tcolorbox}%
}

% Inline kód — underscore package zajistí, ze _ funguje v texttt
\usepackage{underscore}
\newcommand{\pyinline}[1]{\texttt{#1}}
% Pojem (zvýraznění termínu) — stejné jako v preamble-prez.tex
\newcommand{\pojem}[1]{\textbf{#1}}

% Místo pro odpověď studenta
\newcommand{\odpoved}[1][3cm]{%
  \vspace{#1}%
}

% --- Záhlaví / zápatí ---
\usepackage{fancyhdr}
\pagestyle{fancy}
\fancyhf{}
\renewcommand{\headrulewidth}{0.4pt}
\fancyhead[L]{\small\textcolor{microbit-blue}{\textbf{Úvod do Pythonu na Micro:bitu}}}
\fancyhead[R]{\small\textcolor{gray}{\lekce}}
\fancyfoot[C]{\small\thepage}

% --- Ostatní ---
\usepackage{enumitem}
\usepackage{graphicx}
\usepackage{hyperref}
\hypersetup{
  colorlinks=true,
  linkcolor=microbit-blue,
  urlcolor=microbit-blue,
}


\documentclass[a4paper,12pt]{article}

% --- Jazyk a font ---
\usepackage{polyglossia}
\setmainlanguage{czech}
\usepackage{fontspec}

% --- Okraje ---
\usepackage[top=20mm, bottom=20mm, left=22mm, right=22mm]{geometry}

% --- Česky korektní uvozovky ---
\usepackage{csquotes}
\newcommand{\uv}[1]{\enquote{#1}}

% --- Typografie ---
\usepackage{microtype}
\usepackage{parskip}          % odstavce odděleny mezerou, ne odsazením

% --- Barvy a grafika ---
\usepackage[table]{xcolor}
\definecolor{microbit-blue}{HTML}{1E4D8C}
\definecolor{code-bg}{HTML}{F5F5F5}
\definecolor{task-bg}{HTML}{EEF4FF}

% --- Tabulky ---
\usepackage{booktabs}
\usepackage{multirow}

% --- Zdrojový kód (Python) ---
\usepackage{listings}
\lstset{
  language=Python,
  basicstyle=\ttfamily\small,
  backgroundcolor=\color{code-bg},
  frame=single,
  framerule=0pt,
  rulecolor=\color{gray!30},
  xleftmargin=6pt,
  xrightmargin=6pt,
  breaklines=true,
  showstringspaces=false,
  keywordstyle=\color{microbit-blue}\bfseries,
  stringstyle=\color{teal},
  commentstyle=\color{gray}\itshape,
  literate=
    {á}{{\'a}}1 {é}{{\'e}}1 {í}{{\'i}}1 {ó}{{\'o}}1 {ú}{{\'u}}1
    {ů}{{\r{u}}}1 {č}{{\v{c}}}1 {š}{{\v{s}}}1 {ž}{{\v{z}}}1
    {Á}{{\'A}}1 {É}{{\'E}}1 {Í}{{\'I}}1 {Ó}{{\'O}}1 {Ú}{{\'U}}1
    {Č}{{\v{C}}}1 {Š}{{\v{S}}}1 {Ž}{{\v{Z}}}1 {ň}{{\v{n}}}1 {ř}{{\v{r}}}1,
}

% --- Zadání úkolů ---
\usepackage[most]{tcolorbox}
\tcbuselibrary{breakable}

\newcounter{ukol}[section]
\newenvironment{ukol}[1][]{%
  \stepcounter{ukol}%
  \begin{tcolorbox}[
    colback=task-bg,
    colframe=microbit-blue,
    fonttitle=\bfseries,
    title={Úkol~\theukol\ifx\relax#1\relax\else{: #1}\fi},
    breakable,
    left=6pt, right=6pt, top=4pt, bottom=4pt
  ]%
}{%
  \end{tcolorbox}%
}

% Inline kód — underscore package zajistí, ze _ funguje v texttt
\usepackage{underscore}
\newcommand{\pyinline}[1]{\texttt{#1}}
% Pojem (zvýraznění termínu) — stejné jako v preamble-prez.tex
\newcommand{\pojem}[1]{\textbf{#1}}

% Místo pro odpověď studenta
\newcommand{\odpoved}[1][3cm]{%
  \vspace{#1}%
}

% --- Záhlaví / zápatí ---
\usepackage{fancyhdr}
\pagestyle{fancy}
\fancyhf{}
\renewcommand{\headrulewidth}{0.4pt}
\fancyhead[L]{\small\textcolor{microbit-blue}{\textbf{Úvod do Pythonu na Micro:bitu}}}
\fancyhead[R]{\small\textcolor{gray}{\lekce}}
\fancyfoot[C]{\small\thepage}

% --- Ostatní ---
\usepackage{enumitem}
\usepackage{graphicx}
\usepackage{hyperref}
\hypersetup{
  colorlinks=true,
  linkcolor=microbit-blue,
  urlcolor=microbit-blue,
}


\documentclass[a4paper,12pt]{article}

% --- Jazyk a font ---
\usepackage{polyglossia}
\setmainlanguage{czech}
\usepackage{fontspec}

% --- Okraje ---
\usepackage[top=20mm, bottom=20mm, left=22mm, right=22mm]{geometry}

% --- Česky korektní uvozovky ---
\usepackage{csquotes}
\newcommand{\uv}[1]{\enquote{#1}}

% --- Typografie ---
\usepackage{microtype}
\usepackage{parskip}          % odstavce odděleny mezerou, ne odsazením

% --- Barvy a grafika ---
\usepackage[table]{xcolor}
\definecolor{microbit-blue}{HTML}{1E4D8C}
\definecolor{code-bg}{HTML}{F5F5F5}
\definecolor{task-bg}{HTML}{EEF4FF}

% --- Tabulky ---
\usepackage{booktabs}
\usepackage{multirow}

% --- Zdrojový kód (Python) ---
\usepackage{listings}
\lstset{
  language=Python,
  basicstyle=\ttfamily\small,
  backgroundcolor=\color{code-bg},
  frame=single,
  framerule=0pt,
  rulecolor=\color{gray!30},
  xleftmargin=6pt,
  xrightmargin=6pt,
  breaklines=true,
  showstringspaces=false,
  keywordstyle=\color{microbit-blue}\bfseries,
  stringstyle=\color{teal},
  commentstyle=\color{gray}\itshape,
  literate=
    {á}{{\'a}}1 {é}{{\'e}}1 {í}{{\'i}}1 {ó}{{\'o}}1 {ú}{{\'u}}1
    {ů}{{\r{u}}}1 {č}{{\v{c}}}1 {š}{{\v{s}}}1 {ž}{{\v{z}}}1
    {Á}{{\'A}}1 {É}{{\'E}}1 {Í}{{\'I}}1 {Ó}{{\'O}}1 {Ú}{{\'U}}1
    {Č}{{\v{C}}}1 {Š}{{\v{S}}}1 {Ž}{{\v{Z}}}1 {ň}{{\v{n}}}1 {ř}{{\v{r}}}1,
}

% --- Zadání úkolů ---
\usepackage[most]{tcolorbox}
\tcbuselibrary{breakable}

\newcounter{ukol}[section]
\newenvironment{ukol}[1][]{%
  \stepcounter{ukol}%
  \begin{tcolorbox}[
    colback=task-bg,
    colframe=microbit-blue,
    fonttitle=\bfseries,
    title={Úkol~\theukol\ifx\relax#1\relax\else{: #1}\fi},
    breakable,
    left=6pt, right=6pt, top=4pt, bottom=4pt
  ]%
}{%
  \end{tcolorbox}%
}

% Inline kód — underscore package zajistí, ze _ funguje v texttt
\usepackage{underscore}
\newcommand{\pyinline}[1]{\texttt{#1}}
% Pojem (zvýraznění termínu) — stejné jako v preamble-prez.tex
\newcommand{\pojem}[1]{\textbf{#1}}

% Místo pro odpověď studenta
\newcommand{\odpoved}[1][3cm]{%
  \vspace{#1}%
}

% --- Záhlaví / zápatí ---
\usepackage{fancyhdr}
\pagestyle{fancy}
\fancyhf{}
\renewcommand{\headrulewidth}{0.4pt}
\fancyhead[L]{\small\textcolor{microbit-blue}{\textbf{Úvod do Pythonu na Micro:bitu}}}
\fancyhead[R]{\small\textcolor{gray}{\lekce}}
\fancyfoot[C]{\small\thepage}

% --- Ostatní ---
\usepackage{enumitem}
\usepackage{graphicx}
\usepackage{hyperref}
\hypersetup{
  colorlinks=true,
  linkcolor=microbit-blue,
  urlcolor=microbit-blue,
}


\begin{document}

\begin{center}
  {\Large\bfseries\textcolor{microbit-blue}{Lekce 04: Cykly a animace}}\\[4pt]
  {\large Pracovní list}
\end{center}

\medskip\noindent
Cykly umožňují opakovat části programu — buď donekonečna, nebo přesně
zadaný počet opakování. Dnes je využijeme k~animacím na~displeji.

% ------------------------------------------------------------------
\begin{ukol}[Blikající srdce]

Zadej program a spusť ho:

\begin{lstlisting}
from microbit import *

while True:
    display.show(Image.HEART)
    sleep(500)
    display.clear()
    sleep(500)
\end{lstlisting}

\begin{enumerate}
  \item Co dělá \pyinline{sleep(500)}? Co se stane, když číslo změníš na \pyinline{100}?
        Na \pyinline{2000}?
  \item Co dělá \pyinline{display.clear()}?
  \item Upravte program tak, aby srdce blikalo rychle třikrát, pak se na sekundu
        zastavilo — a tak stále dokola. (Tip: vnořený \pyinline{for} cyklus.)
\end{enumerate}

\odpoved[2cm]

\begin{tcolorbox}[colback=code-bg, colframe=gray!40,
                  fonttitle=\small\bfseries, title=sleep() a display.clear()]
  \small
  \begin{itemize}[leftmargin=*]
    \item \pyinline{sleep(ms)} — program čeká zadaný počet milisekund.
          \pyinline{sleep(1000)} = 1 sekunda.
    \item \pyinline{display.clear()} — zhasne všechny LED.
  \end{itemize}
\end{tcolorbox}

\end{ukol}

% ------------------------------------------------------------------
\begin{ukol}[Cyklus for a range()]

\pyinline{for} cyklus opakuje blok přesně tolikrát, kolik řekneme:

\begin{lstlisting}
from microbit import *

for i in range(5):
    display.show(str(i))
    sleep(600)

display.show(Image.YES)
\end{lstlisting}

\begin{enumerate}
  \item Spusť program. Jaká čísla se zobrazí? (Pozor — začínají od 0!)
  \item Změň \pyinline{range(5)} na \pyinline{range(10)} — co se změní?
  \item Jaký je rozdíl mezi \pyinline{while True:} a \pyinline{for i in range(5):}?
\end{enumerate}

\odpoved[2.5cm]

\begin{tcolorbox}[colback=code-bg, colframe=gray!40,
                  fonttitle=\small\bfseries, title=range() a proměnná i]
  \small
  \pyinline{range(n)} vygeneruje čísla \pyinline{0, 1, 2, ..., n-1}.\\
  Proměnná \pyinline{i} (nebo jakékoli jiné jméno) přijímá postupně každou hodnotu.\\
  \pyinline{for i in range(3):} → i bude 0, pak 1, pak 2.
\end{tcolorbox}

\end{ukol}

% ------------------------------------------------------------------
\begin{ukol}[Vlastní animace]

Vytvoř animaci ze sekvence obrázků. Obrázky ulož do~proměnné jako seznam
a~procházej je \pyinline{for} cyklem:

\begin{lstlisting}
from microbit import *

snimky = [Image.HAPPY, Image.SMILE, Image.SAD]

while True:
    for obrazek in snimky:
        display.show(obrazek)
        sleep(400)
\end{lstlisting}

\begin{enumerate}
  \item Spusť program a pozoruj animaci.
  \item Přidej do seznamu \pyinline{snimky} další obrázky dle vlastního výběru.
        Zkus např.: \pyinline{Image.SURPRISED}, \pyinline{Image.ANGRY},
        \pyinline{Image.ASLEEP}, \pyinline{Image.GHOST}, \pyinline{Image.PACMAN}.
  \item Uprav \pyinline{sleep()} tak, aby animace vypadala přirozeně.
\end{enumerate}

\medskip
Výslednou animaci ulož a připrav k odevzdání do Google Classroom.

\end{ukol}

% ------------------------------------------------------------------
\begin{ukol}[Bonus: pohybující se bod]

Funkce \pyinline{display.set_pixel(x, y, jas)} rozsvítí jednu LED na souřadnicích
\pyinline{x} (sloupec 0--4, zleva) a \pyinline{y} (řádek 0--4, shora) s jasem 0--9.

Zkus napsat program, který pohybuje jasným bodem zleva doprava po prostřední řadě:

\begin{lstlisting}
from microbit import *

while True:
    for x in range(5):
        display.clear()
        display.set_pixel(x, 2, 9)
        sleep(150)
\end{lstlisting}

Jak bys program upravil, aby se bod pohyboval tam a zpět?

\odpoved[2.5cm]

\end{ukol}

\end{document}
